\section{链接、运行时与加载：从符号图到进程映像}
\label{sec:linker-loader}

链接器面对的输入不是源语言语句，而是若干节、符号、重定位与归档索引。它所拥有的信息比单个编译单元更全：所有选中对象的符号图和近最终布局已经可见；但它通常又不知道 LLVM IR 中的高级类型、循环语义或 \texttt{nsw} 假设。因此，链接是“全映像推理”，不天然等于全程序优化。 经典链接器文献对符号、重定位和装载的抽象仍有解释力\cite{bin-levine-linkers}，但精确的 RISC-V 与现代 Linux 行为必须回到当前 psABI、gABI 和工具文档。本节把链接时（link time）与装载时（load time）严格分开，并解释 C 运行时如何把 ELF 入口桥接到 \texttt{main}。

\subsection{链接器工作的依赖顺序}

一次静态链接可抽象为如下相互依赖的步骤：

\begin{enumerate}
  \item 读取目标格式与架构属性，拒绝不兼容的机器、浮点 ABI 或对象类别；
  \item 解析全局/弱符号并按需抽取归档成员；
  \item 若启用节垃圾回收（section garbage collection），从入口和保留根标记活跃输入节；
  \item 依据链接脚本把输入节归并到输出节，并分配文件偏移和虚拟地址；
  \item 合成 GOT、PLT、动态表或边界符号等结构；
  \item 计算重定位并执行目标相关松弛；若代码缩短，重新收敛受影响布局；
  \item 写出 ELF 头、程序头、输出节与可选链接映射文件。
\end{enumerate}

真实链接器会交错其中若干步骤，但依赖不可颠倒：符号解析影响归档成员和活跃节，布局决定 \(S\) 与 \(P\)，松弛又会改变后继地址。GNU \texttt{ld} 总会使用链接脚本；未给 \texttt{-T} 时采用可由 \texttt{--verbose} 查看的一套内建脚本\cite{bin-gnu-ld-scripts}。在托管 Linux 实验中应让 \texttt{riscv64-linux-gnu-gcc} 驱动最终链接，因为驱动同时选择正确的启动对象、libgcc、libc、ABI、sysroot 与默认脚本。直接调用 \texttt{ld} 只看似“更底层”，实际很容易漏掉语言运行时契约。

预检完整模式使用

\begin{verbatim}
riscv64-linux-gnu-gcc -static -no-pie -march=rv64gc -mabi=lp64d \
  handwritten-asm.rv64.o libsysy-glibc-rv64.a \
  -Wl,-Map=asm.link.map \
  -o clipped-dot-asm.rv64
\end{verbatim}

链接映射（link map）把“哪个输入成员因何进入哪个输出节、最终地址是多少”保存为可审计证据；仅看最终反汇编无法可靠恢复归档抽取决策。\texttt{-static -no-pie} 是本次实验为了减少装载期变量而选定的观察条件，不是 RISC-V 程序的普遍要求。

\subsection{符号解析与课程运行时的实际不兼容}

\observed 对课程提供的 \texttt{lib/libsysy\_riscv.a} 运行 \texttt{nm -u}，其唯一成员 \texttt{sylib.o} 除 libc 接口外还含未定义的 \texttt{\_impure\_ptr}；同一成员确实定义 \texttt{getint}、\texttt{putint} 和 \texttt{putch}。\texttt{\_impure\_ptr} 是 newlib 可重入状态机制的一部分，而 Ubuntu/Debian 的 \texttt{riscv64-linux-gnu-gcc} 交叉 sysroot 使用 glibc。架构同为 RISC-V、函数名也相同，并不保证 C 库 ABI 与内部符号生态兼容；直接把这个 newlib 对象交给 glibc 链接环境，会留下无法满足的 \texttt{\_impure\_ptr}。

正确修复不是伪造一个同名全局变量，也不是忽略未解析引用。预检从仓库中同一份 \texttt{lib/sylib.c} 用当前 \texttt{rv64gc/lp64d} Linux 交叉工具链重新编译为 \texttt{sylib.glibc.rv64.o}，再以 \texttt{ar rcs} 构成 \texttt{libsysy-glibc-rv64.a}。\observed 对重建归档执行 \texttt{nm}，\texttt{getint}/\texttt{putint}/\texttt{putch} 仍为强文本定义，而不再出现 \texttt{\_impure\_ptr}。这保留了公开 SysY 接口，却使实现的 libc 依赖与最终 sysroot 一致。

这个案例揭示了二进制兼容的三层检查：

\begin{table}[t]
\caption{运行时库兼容性不能由文件名推出}
\label{tab:runtime-compat}
\centering
\small
\begin{tabular}{p{.20\linewidth}p{.33\linewidth}p{.37\linewidth}}
\toprule
层次 & 可检查证据 & 本案例结论\\
\midrule
ISA/ELF & \texttt{readelf -h -A}、\texttt{-march} & 都是 ELF64 RISC-V；仍不足以证明可链接\\
psABI & ELF flags、LP64D、栈与调用规则 & 应保持 \texttt{rv64gc/lp64d} 一致\\
系统 C 库 ABI & \texttt{nm -u}、sysroot 与实际链接 & 提供库依赖 newlib 的 \texttt{\_impure\_ptr}，Debian/Ubuntu glibc 不提供；必须从源码重建\\
\bottomrule
\end{tabular}
\end{table}

重建也带来不能隐藏的行为。\texttt{sylib.c} 以构造器（constructor）初始化计时状态，并以析构器（destructor）在程序终止时输出汇总。即使案例从未调用计时 API，\texttt{after\_main} 仍向标准错误精确写出 \texttt{TOTAL: 0H-0M-0S-0us\textbackslash n}。这不是 \texttt{clipped\_dot} 的标准输出，而是已链接运行时的另一个可观察通道。验证器分别断言 stdout 的六个 oracle 与 stderr 的运行时行，不能为得到“干净结果”而静默丢弃 stderr。

\subsection{布局、重定位与 RISC-V 松弛}

链接布局把每个选中输入节放入输出节，再由程序头把可装载内容分组。常见脚本把 \texttt{.text.*} 合入可执行代码，把只读常量放入只读区域，把 \texttt{.data.*}/\texttt{.bss.*} 放入可写区域；但实际边界、对齐、TLS、构造器数组与 unwind 信息由平台脚本共同决定。一个只含四个节的裸机脚本很适合教学，却不是 Linux 生产脚本的替代品：它通常缺少 \texttt{crt} 启动、动态元数据、TLS、\texttt{.init\_array}/\texttt{.fini\_array} 与合适的程序头。

重定位应用后，链接器已知 \texttt{getint} 等定义的最终地址。若某个 \texttt{AUIPC+JALR} 调用落入 JAL 可达范围，并且对象以 \texttt{R\_RISCV\_RELAX} 授权，链接器可以把两条指令改为一条 \texttt{jal}。这称为链接器松弛（linker relaxation）：前端为一般距离生成保守序列，链接器凭最终布局消除不再需要的通用性\cite{bin-riscv-psabi-relax}。松弛可能删除字节，故后方符号地址和对齐填充也需更新；相关重定位必须成组一致变换。

这类优化的证据应比较同一对象在正常链接与 \texttt{--no-relax} 下的最终反汇编和大小，而不能只根据源汇编的 \texttt{call} 猜测。也不应许诺固定节省比例：是否可缩短取决于距离、链接器版本、C 扩展、代码模型和布局。链接器能证明的是地址范围，不是被调函数语义；跨函数常量传播和内联属于链接时优化（link-time optimization, LTO）在 IR 层承担的另一问题。

\subsection{最终 ELF：节服务链接，段服务装载}

最终静态可执行文件通常是 \texttt{ET\_EXEC}，位置无关可执行文件常是 \texttt{ET\_DYN}。无论哪种，加载器主要读取程序头。对一个 \texttt{PT\_LOAD}，文件区间 \([p_{offset},p_{offset}+p_{filesz})\) 提供初始字节，内存区间长度为 \(p_{memsz}\)；若 \(p_{memsz}>p_{filesz}\)，尾部由系统清零。\texttt{p\_vaddr} 与 \texttt{p\_offset} 必须按规定对齐同余\cite{bin-sysv-gabi}。于是 \texttt{.bss} 的准确解释是：链接视图中它通常是 \texttt{SHT\_NOBITS}，装载视图中其空间往往表现为可写 \texttt{PT\_LOAD} 的零填充尾部。

\begin{table}[t]
\caption{两种视图与两个时间点}
\label{tab:section-segment}
\centering
\small
\begin{tabular}{p{.19\linewidth}p{.33\linewidth}p{.37\linewidth}}
\toprule
概念 & 主要消费者 & 回答的问题\\
\midrule
节（section） & 汇编器、静态链接器、调试/分析工具 & 哪些字节属于代码、符号、重定位或调试信息？\\
段（segment） & 内核 ELF 加载器、动态链接器 & 哪些文件页映到哪些虚拟地址，具有什么权限？\\
链接时 & 静态链接器 & 哪个定义满足引用，输出地址与静态重定位为何？\\
装载时 & 内核与用户态动态链接器 & 本次执行的装载基址、共享对象与动态重定位为何？\\
\bottomrule
\end{tabular}
\end{table}

关键程序头还包括：\texttt{PT\_INTERP} 给出动态解释器路径，\texttt{PT\_DYNAMIC} 指向动态链接元数据，\texttt{PT\_TLS} 描述线程局部存储模板；Linux/GNU 扩展 \texttt{PT\_GNU\_STACK} 表达栈可执行性请求，\texttt{PT\_GNU\_RELRO} 标记重定位后应改为只读的区域。\texttt{readelf -l} 的 Section-to-Segment mapping 可以展示节如何被成组覆盖，但操作系统并不需要按节名逐一装载。

\subsection{从 \texttt{execve} 到 \texttt{main}}

成功的 \texttt{execve} 不会创建新进程，而是以新程序映像替换当前进程映像，保留的 PID 与许多内核对象遵循 Linux 明确定义的规则\cite{bin-linux-execve}。内核验证 ELF，选择装载基址并映射 \texttt{PT\_LOAD}，建立初始栈（参数、环境和辅助向量），然后把控制权交给 ELF 入口或解释器。

ELF 头的 \texttt{e\_entry} 通常指向 \texttt{\_start}，不是 C 的 \texttt{main}。由编译器驱动加入的 C 运行时启动对象（C runtime startup, CRT）接收初始机器状态，建立 libc 环境，运行初始化数组，间接调用 \texttt{main}，再把返回值转换为进程终止并运行析构序列。本文重建运行时中的 \texttt{before\_main}/\texttt{after\_main} 正是经 \texttt{.init\_array}/\texttt{.fini\_array} 参与该链条；若直接把 ELF 入口设为 \texttt{main}，这些契约不会自动发生。

静态与动态链接的差异可沿责任边界理解：

\begin{itemize}
  \item 完全静态可执行文件已经在构建时纳入所需库代码并完成普通外部符号解析，通常没有 \texttt{PT\_INTERP}；文件较大，但部署时不依赖相应共享库解析器。它仍由内核装载，仍有 CRT，并不意味着“没有运行时”。
  \item 动态可执行文件保留 \texttt{PT\_INTERP}、\texttt{PT\_DYNAMIC}、动态符号与动态重定位。内核映射程序与解释器后，用户态动态链接器加载 \texttt{DT\_NEEDED} 共享对象，解析符号和重定位，运行初始化，再转交程序入口\cite{bin-sysv-gabi-dynamic}。
  \item PIE 通常以 \texttt{ET\_DYN} 表示，实际虚拟地址为链接时相对值加本次装载基址 \(B\)。例如 \texttt{R\_RISCV\_RELATIVE} 可按 \(B+A\) 应用而无需名字查找。
\end{itemize}

因此“链接器执行机器码生成可执行文件”和“加载器调用 \texttt{main}”都不准确。链接器变换并写出文件；内核与动态链接器建立映像；CRT 才建立语言级入口条件。

\subsection{验证闭环与证据强度}

完整预检对 C 观察载体、以 C 前端承载的 SysY 源、手写 LLVM IR 和手写汇编四条路径分别生成对象、链接映射、静态 RV64GC 可执行文件和最终反汇编，再用 QEMU Linux-user 执行。Clang 生成的 \code{.s} 在本实验中用于比较观察，并未作为第五种独立可执行形式。输入 \(-1,0,1,4,8,10\) 的预期 stdout 分别为 \(0,0,-8,8,16,16\)，且退出码为零；stderr 另有上述计时汇总。QEMU 验证的是用户态可观察行为，不替代 ELF 静态结构检查；反过来，\texttt{readelf} 证明结构合法，也不能证明算法输出正确。把二者并置才形成“结构契约 + 行为 oracle”的闭环。

证据仍有边界：六个样例不是对所有输入的形式化等价证明，QEMU 也不是特定物理微架构的性能证据。更强的后续工作可用差分测试随机生成 \(n\)，用独立参考模型计算结果；对限定输入域 \([0,8]\) 甚至可穷举。若评估松弛或 C 扩展的性能，则需固定链接器、处理器/模拟器、测量方法与置信区间，避免把代码大小改善直接外推为运行时间改善。

最终，一条贯穿本案例的原则是“逐步提交信息”：汇编器提交局部编码并保留重定位；静态链接器提交全映像地址并保留必要的动态元数据；加载器为一次执行提交装载基址与内存映射。每一阶段都不应猜测下一阶段才能知道的事实，也不应丢弃下一阶段证明正确性所需的证据。
